% paper_figures.tex
% Auto-generated by figures_tables.py — RIS M-FSK Sensing Suite
% Include in your main .tex file with:
%   % paper_figures.tex
% Auto-generated by figures_tables.py — RIS M-FSK Sensing Suite
% Include in your main .tex file with:
%   % paper_figures.tex
% Auto-generated by figures_tables.py — RIS M-FSK Sensing Suite
% Include in your main .tex file with:
%   % paper_figures.tex
% Auto-generated by figures_tables.py — RIS M-FSK Sensing Suite
% Include in your main .tex file with:
%   \input{paper_figures}
%
% Requires: \usepackage{graphicx}

% ─── Phase 1 figures ─────────────────────────────────────────────────

\begin{figure}[!t]
  \centering
  \includegraphics[width=\columnwidth]{figures/F4_transmission_amplitude}
  \caption{T-RIS PIN-diode transmission amplitude $|T_n|$ vs frequency. On-state ($V_f > 0$) and off-state ($V_r < 0$) responses at 28 GHz carrier (vertical dashed line).}
  \label{fig:pin_transmission}
\end{figure}

\begin{figure}[!t]
  \centering
  \includegraphics[width=\columnwidth]{figures/F5_phase_response}
  \caption{T-RIS PIN-diode phase response $\angle T_n(f)$ in the on and off states. Phase difference $\Delta\phi = 64.5^\circ$ at 28 GHz.}
  \label{fig:pin_phase}
\end{figure}

\begin{figure}[!t]
  \centering
  \includegraphics[width=\columnwidth]{figures/F6_impedance_magnitude}
  \caption{PIN-diode impedance magnitude $|Z_n(f)|$ for forward and reverse bias.}
  \label{fig:pin_impedance}
\end{figure}

% ─── Phase 2 figures ─────────────────────────────────────────────────

\begin{figure}[!t]
  \centering
  \includegraphics[width=\columnwidth]{figures/F7_mfsk_spectrum}
  \caption{M-FSK tone spectrum at T-RIS output. Twelve sidebands at $f_c + k\Delta f$, $k=1,\ldots,12$, within the BPF passband.}
  \label{fig:mfsk_spectrum}
\end{figure}

\begin{figure*}[!t]
  \centering
  \includegraphics[width=\textwidth]{figures/F8_mfsk_spectrogram}
  \caption{Time--frequency spectrogram of the T-RIS M-FSK waveform. Each row shows one symbol period $T_{\rm sym} = 1.40\ \mu$s.}
  \label{fig:mfsk_spectrogram}
\end{figure*}

% ─── Phase 3 figures ─────────────────────────────────────────────────

\begin{figure}[!t]
  \centering
  \includegraphics[width=\columnwidth]{figures/F10_link_budget_waterfall}
  \caption{Bistatic link budget waterfall. Per-stage contributions (bars) and cumulative signal power (dashed) across the four-hop Tx $\to$ T-RIS $\to$ target $\to$ R-RIS $\to$ Rx path.}
  \label{fig:link_budget}
\end{figure}

\begin{figure*}[!t]
  \centering
  \includegraphics[width=\textwidth]{figures/F11_range_doppler}
  \caption{(a) IFFT range profile from $N_{\rm sw}=1024$ sweeps: target peak at 9.99 m (true 10 m, error $<1$ cm). (b) Range--Doppler map: target at $(10.0\ {\rm m},\ 2.5\ {\rm m/s})$.}
  \label{fig:range_doppler}
\end{figure*}

% ─── Phase 4 figures ─────────────────────────────────────────────────

\begin{figure*}[!t]
  \centering
  \includegraphics[width=\textwidth]{figures/F12_rmse_range_position}
  \caption{(a) Range RMSE and CRB vs single-symbol SNR. (b) 2-D position RMSE and CRB. Vertical dotted line: operating point SNR$_1 = -14.5$ dB.}
  \label{fig:rmse_range_pos}
\end{figure*}

\begin{figure}[!t]
  \centering
  \includegraphics[width=\columnwidth]{figures/F13_rmse_velocity}
  \caption{Velocity RMSE, CRB, and FFT estimation floor vs SNR. Three regimes: noise-dominated, CRB-tracking, and FFT-bin-limited.}
  \label{fig:rmse_velocity}
\end{figure}

% ─── Phase 5 figures ─────────────────────────────────────────────────

\begin{figure*}[!t]
  \centering
  \includegraphics[width=\textwidth]{figures/F14_tracking_trajectory}
  \caption{(a) Closed-loop beam-tracking scene over 50 CPIs ($T_{\rm tot}=860$ ms). Beam-steering arrows and colour-coded estimates. (b) Angle commands: true $\theta_{\rm true}$, noisy $\hat{\theta}$, and $\alpha$-smoothed $\tilde{\theta}$ ($\alpha=0.3$) vs CPI frame.}
  \label{fig:tracking_traj}
\end{figure*}

\begin{figure*}[!t]
  \centering
  \includegraphics[width=\textwidth]{figures/F15_tracking_errors}
  \caption{(a) Position-error scatter: 50th-percentile circle $= 5.3$ cm, 95th $= 9.2$ cm. (b) Empirical CDFs of position, range, and tangential errors; CRB reference $\sigma_R = 1.32$ cm.}
  \label{fig:tracking_err}
\end{figure*}

% ─── Phase 6 figures ─────────────────────────────────────────────────

\begin{figure*}[!t]
  \centering
  \includegraphics[width=\textwidth]{figures/F16_multitarget_vs_fmcw}
  \caption{Multi-target range profiles. (a) M-FSK IFFT: three clean peaks at true ranges 4, 7, 10 m regardless of target velocity. (b) FMCW FFT (same $B$, same CPI): Doppler-induced displacements of $+4.82$ m and $-3.21$ m create ghost targets.}
  \label{fig:multitarget}
\end{figure*}

\begin{figure}[!t]
  \centering
  \includegraphics[width=\columnwidth]{figures/F17_power_comparison}
  \caption{Stacked power breakdown for four sensing architectures. The two-stage T-RIS/R-RIS system achieves $2.33$ W, a $65\%$ reduction vs AESA ($6.56$ W).}
  \label{fig:power_bar}
\end{figure}

\begin{figure}[!t]
  \centering
  \includegraphics[width=\columnwidth]{figures/F18_spider_chart}
  \caption{Normalised performance spider chart (five axes). The proposed two-stage architecture (solid blue) leads in symbol rate and beam agility while matching the baseline in range resolution and hardware cost.}
  \label{fig:spider}
\end{figure}

% ─── Phase 7 figures ─────────────────────────────────────────────────

\begin{figure*}[!t]
  \centering
  \includegraphics[width=\textwidth]{figures/FS1_element_tolerance}
  \caption{Element-level manufacturing tolerance robustness. (a) Beam-gain change vs amplitude error $\sigma_A$: negligible impact ($<0.001$ dB at $\sigma_A=2\%$). (b) Beam-gain loss vs phase error $\sigma_\phi$: $0.012$ dB at nominal $3^\circ$ tolerance.}
  \label{fig:tolerance}
\end{figure*}

\begin{figure*}[!t]
  \centering
  \includegraphics[width=\textwidth]{figures/FS2_gap_and_quantisation}
  \caption{Hardware sensitivity. (a) SNR$_{\rm CPI}$ and CRB$_R$ vs inter-stage gap $d_g$: $5.10$ dB/mm sensitivity drives $\pm 0.196$ mm assembly tolerance. (b) Beam-gain loss vs phase bits $B$: $B=3$ bits incurs only $0.22$ dB loss.}
  \label{fig:gap_quant}
\end{figure*}

%
% Requires: \usepackage{graphicx}

% ─── Phase 1 figures ─────────────────────────────────────────────────

\begin{figure}[!t]
  \centering
  \includegraphics[width=\columnwidth]{figures/F4_transmission_amplitude}
  \caption{T-RIS PIN-diode transmission amplitude $|T_n|$ vs frequency. On-state ($V_f > 0$) and off-state ($V_r < 0$) responses at 28 GHz carrier (vertical dashed line).}
  \label{fig:pin_transmission}
\end{figure}

\begin{figure}[!t]
  \centering
  \includegraphics[width=\columnwidth]{figures/F5_phase_response}
  \caption{T-RIS PIN-diode phase response $\angle T_n(f)$ in the on and off states. Phase difference $\Delta\phi = 64.5^\circ$ at 28 GHz.}
  \label{fig:pin_phase}
\end{figure}

\begin{figure}[!t]
  \centering
  \includegraphics[width=\columnwidth]{figures/F6_impedance_magnitude}
  \caption{PIN-diode impedance magnitude $|Z_n(f)|$ for forward and reverse bias.}
  \label{fig:pin_impedance}
\end{figure}

% ─── Phase 2 figures ─────────────────────────────────────────────────

\begin{figure}[!t]
  \centering
  \includegraphics[width=\columnwidth]{figures/F7_mfsk_spectrum}
  \caption{M-FSK tone spectrum at T-RIS output. Twelve sidebands at $f_c + k\Delta f$, $k=1,\ldots,12$, within the BPF passband.}
  \label{fig:mfsk_spectrum}
\end{figure}

\begin{figure*}[!t]
  \centering
  \includegraphics[width=\textwidth]{figures/F8_mfsk_spectrogram}
  \caption{Time--frequency spectrogram of the T-RIS M-FSK waveform. Each row shows one symbol period $T_{\rm sym} = 1.40\ \mu$s.}
  \label{fig:mfsk_spectrogram}
\end{figure*}

% ─── Phase 3 figures ─────────────────────────────────────────────────

\begin{figure}[!t]
  \centering
  \includegraphics[width=\columnwidth]{figures/F10_link_budget_waterfall}
  \caption{Bistatic link budget waterfall. Per-stage contributions (bars) and cumulative signal power (dashed) across the four-hop Tx $\to$ T-RIS $\to$ target $\to$ R-RIS $\to$ Rx path.}
  \label{fig:link_budget}
\end{figure}

\begin{figure*}[!t]
  \centering
  \includegraphics[width=\textwidth]{figures/F11_range_doppler}
  \caption{(a) IFFT range profile from $N_{\rm sw}=1024$ sweeps: target peak at 9.99 m (true 10 m, error $<1$ cm). (b) Range--Doppler map: target at $(10.0\ {\rm m},\ 2.5\ {\rm m/s})$.}
  \label{fig:range_doppler}
\end{figure*}

% ─── Phase 4 figures ─────────────────────────────────────────────────

\begin{figure*}[!t]
  \centering
  \includegraphics[width=\textwidth]{figures/F12_rmse_range_position}
  \caption{(a) Range RMSE and CRB vs single-symbol SNR. (b) 2-D position RMSE and CRB. Vertical dotted line: operating point SNR$_1 = -14.5$ dB.}
  \label{fig:rmse_range_pos}
\end{figure*}

\begin{figure}[!t]
  \centering
  \includegraphics[width=\columnwidth]{figures/F13_rmse_velocity}
  \caption{Velocity RMSE, CRB, and FFT estimation floor vs SNR. Three regimes: noise-dominated, CRB-tracking, and FFT-bin-limited.}
  \label{fig:rmse_velocity}
\end{figure}

% ─── Phase 5 figures ─────────────────────────────────────────────────

\begin{figure*}[!t]
  \centering
  \includegraphics[width=\textwidth]{figures/F14_tracking_trajectory}
  \caption{(a) Closed-loop beam-tracking scene over 50 CPIs ($T_{\rm tot}=860$ ms). Beam-steering arrows and colour-coded estimates. (b) Angle commands: true $\theta_{\rm true}$, noisy $\hat{\theta}$, and $\alpha$-smoothed $\tilde{\theta}$ ($\alpha=0.3$) vs CPI frame.}
  \label{fig:tracking_traj}
\end{figure*}

\begin{figure*}[!t]
  \centering
  \includegraphics[width=\textwidth]{figures/F15_tracking_errors}
  \caption{(a) Position-error scatter: 50th-percentile circle $= 5.3$ cm, 95th $= 9.2$ cm. (b) Empirical CDFs of position, range, and tangential errors; CRB reference $\sigma_R = 1.32$ cm.}
  \label{fig:tracking_err}
\end{figure*}

% ─── Phase 6 figures ─────────────────────────────────────────────────

\begin{figure*}[!t]
  \centering
  \includegraphics[width=\textwidth]{figures/F16_multitarget_vs_fmcw}
  \caption{Multi-target range profiles. (a) M-FSK IFFT: three clean peaks at true ranges 4, 7, 10 m regardless of target velocity. (b) FMCW FFT (same $B$, same CPI): Doppler-induced displacements of $+4.82$ m and $-3.21$ m create ghost targets.}
  \label{fig:multitarget}
\end{figure*}

\begin{figure}[!t]
  \centering
  \includegraphics[width=\columnwidth]{figures/F17_power_comparison}
  \caption{Stacked power breakdown for four sensing architectures. The two-stage T-RIS/R-RIS system achieves $2.33$ W, a $65\%$ reduction vs AESA ($6.56$ W).}
  \label{fig:power_bar}
\end{figure}

\begin{figure}[!t]
  \centering
  \includegraphics[width=\columnwidth]{figures/F18_spider_chart}
  \caption{Normalised performance spider chart (five axes). The proposed two-stage architecture (solid blue) leads in symbol rate and beam agility while matching the baseline in range resolution and hardware cost.}
  \label{fig:spider}
\end{figure}

% ─── Phase 7 figures ─────────────────────────────────────────────────

\begin{figure*}[!t]
  \centering
  \includegraphics[width=\textwidth]{figures/FS1_element_tolerance}
  \caption{Element-level manufacturing tolerance robustness. (a) Beam-gain change vs amplitude error $\sigma_A$: negligible impact ($<0.001$ dB at $\sigma_A=2\%$). (b) Beam-gain loss vs phase error $\sigma_\phi$: $0.012$ dB at nominal $3^\circ$ tolerance.}
  \label{fig:tolerance}
\end{figure*}

\begin{figure*}[!t]
  \centering
  \includegraphics[width=\textwidth]{figures/FS2_gap_and_quantisation}
  \caption{Hardware sensitivity. (a) SNR$_{\rm CPI}$ and CRB$_R$ vs inter-stage gap $d_g$: $5.10$ dB/mm sensitivity drives $\pm 0.196$ mm assembly tolerance. (b) Beam-gain loss vs phase bits $B$: $B=3$ bits incurs only $0.22$ dB loss.}
  \label{fig:gap_quant}
\end{figure*}

%
% Requires: \usepackage{graphicx}

% ─── Phase 1 figures ─────────────────────────────────────────────────

\begin{figure}[!t]
  \centering
  \includegraphics[width=\columnwidth]{figures/F4_transmission_amplitude}
  \caption{T-RIS PIN-diode transmission amplitude $|T_n|$ vs frequency. On-state ($V_f > 0$) and off-state ($V_r < 0$) responses at 28 GHz carrier (vertical dashed line).}
  \label{fig:pin_transmission}
\end{figure}

\begin{figure}[!t]
  \centering
  \includegraphics[width=\columnwidth]{figures/F5_phase_response}
  \caption{T-RIS PIN-diode phase response $\angle T_n(f)$ in the on and off states. Phase difference $\Delta\phi = 64.5^\circ$ at 28 GHz.}
  \label{fig:pin_phase}
\end{figure}

\begin{figure}[!t]
  \centering
  \includegraphics[width=\columnwidth]{figures/F6_impedance_magnitude}
  \caption{PIN-diode impedance magnitude $|Z_n(f)|$ for forward and reverse bias.}
  \label{fig:pin_impedance}
\end{figure}

% ─── Phase 2 figures ─────────────────────────────────────────────────

\begin{figure}[!t]
  \centering
  \includegraphics[width=\columnwidth]{figures/F7_mfsk_spectrum}
  \caption{M-FSK tone spectrum at T-RIS output. Twelve sidebands at $f_c + k\Delta f$, $k=1,\ldots,12$, within the BPF passband.}
  \label{fig:mfsk_spectrum}
\end{figure}

\begin{figure*}[!t]
  \centering
  \includegraphics[width=\textwidth]{figures/F8_mfsk_spectrogram}
  \caption{Time--frequency spectrogram of the T-RIS M-FSK waveform. Each row shows one symbol period $T_{\rm sym} = 1.40\ \mu$s.}
  \label{fig:mfsk_spectrogram}
\end{figure*}

% ─── Phase 3 figures ─────────────────────────────────────────────────

\begin{figure}[!t]
  \centering
  \includegraphics[width=\columnwidth]{figures/F10_link_budget_waterfall}
  \caption{Bistatic link budget waterfall. Per-stage contributions (bars) and cumulative signal power (dashed) across the four-hop Tx $\to$ T-RIS $\to$ target $\to$ R-RIS $\to$ Rx path.}
  \label{fig:link_budget}
\end{figure}

\begin{figure*}[!t]
  \centering
  \includegraphics[width=\textwidth]{figures/F11_range_doppler}
  \caption{(a) IFFT range profile from $N_{\rm sw}=1024$ sweeps: target peak at 9.99 m (true 10 m, error $<1$ cm). (b) Range--Doppler map: target at $(10.0\ {\rm m},\ 2.5\ {\rm m/s})$.}
  \label{fig:range_doppler}
\end{figure*}

% ─── Phase 4 figures ─────────────────────────────────────────────────

\begin{figure*}[!t]
  \centering
  \includegraphics[width=\textwidth]{figures/F12_rmse_range_position}
  \caption{(a) Range RMSE and CRB vs single-symbol SNR. (b) 2-D position RMSE and CRB. Vertical dotted line: operating point SNR$_1 = -14.5$ dB.}
  \label{fig:rmse_range_pos}
\end{figure*}

\begin{figure}[!t]
  \centering
  \includegraphics[width=\columnwidth]{figures/F13_rmse_velocity}
  \caption{Velocity RMSE, CRB, and FFT estimation floor vs SNR. Three regimes: noise-dominated, CRB-tracking, and FFT-bin-limited.}
  \label{fig:rmse_velocity}
\end{figure}

% ─── Phase 5 figures ─────────────────────────────────────────────────

\begin{figure*}[!t]
  \centering
  \includegraphics[width=\textwidth]{figures/F14_tracking_trajectory}
  \caption{(a) Closed-loop beam-tracking scene over 50 CPIs ($T_{\rm tot}=860$ ms). Beam-steering arrows and colour-coded estimates. (b) Angle commands: true $\theta_{\rm true}$, noisy $\hat{\theta}$, and $\alpha$-smoothed $\tilde{\theta}$ ($\alpha=0.3$) vs CPI frame.}
  \label{fig:tracking_traj}
\end{figure*}

\begin{figure*}[!t]
  \centering
  \includegraphics[width=\textwidth]{figures/F15_tracking_errors}
  \caption{(a) Position-error scatter: 50th-percentile circle $= 5.3$ cm, 95th $= 9.2$ cm. (b) Empirical CDFs of position, range, and tangential errors; CRB reference $\sigma_R = 1.32$ cm.}
  \label{fig:tracking_err}
\end{figure*}

% ─── Phase 6 figures ─────────────────────────────────────────────────

\begin{figure*}[!t]
  \centering
  \includegraphics[width=\textwidth]{figures/F16_multitarget_vs_fmcw}
  \caption{Multi-target range profiles. (a) M-FSK IFFT: three clean peaks at true ranges 4, 7, 10 m regardless of target velocity. (b) FMCW FFT (same $B$, same CPI): Doppler-induced displacements of $+4.82$ m and $-3.21$ m create ghost targets.}
  \label{fig:multitarget}
\end{figure*}

\begin{figure}[!t]
  \centering
  \includegraphics[width=\columnwidth]{figures/F17_power_comparison}
  \caption{Stacked power breakdown for four sensing architectures. The two-stage T-RIS/R-RIS system achieves $2.33$ W, a $65\%$ reduction vs AESA ($6.56$ W).}
  \label{fig:power_bar}
\end{figure}

\begin{figure}[!t]
  \centering
  \includegraphics[width=\columnwidth]{figures/F18_spider_chart}
  \caption{Normalised performance spider chart (five axes). The proposed two-stage architecture (solid blue) leads in symbol rate and beam agility while matching the baseline in range resolution and hardware cost.}
  \label{fig:spider}
\end{figure}

% ─── Phase 7 figures ─────────────────────────────────────────────────

\begin{figure*}[!t]
  \centering
  \includegraphics[width=\textwidth]{figures/FS1_element_tolerance}
  \caption{Element-level manufacturing tolerance robustness. (a) Beam-gain change vs amplitude error $\sigma_A$: negligible impact ($<0.001$ dB at $\sigma_A=2\%$). (b) Beam-gain loss vs phase error $\sigma_\phi$: $0.012$ dB at nominal $3^\circ$ tolerance.}
  \label{fig:tolerance}
\end{figure*}

\begin{figure*}[!t]
  \centering
  \includegraphics[width=\textwidth]{figures/FS2_gap_and_quantisation}
  \caption{Hardware sensitivity. (a) SNR$_{\rm CPI}$ and CRB$_R$ vs inter-stage gap $d_g$: $5.10$ dB/mm sensitivity drives $\pm 0.196$ mm assembly tolerance. (b) Beam-gain loss vs phase bits $B$: $B=3$ bits incurs only $0.22$ dB loss.}
  \label{fig:gap_quant}
\end{figure*}

%
% Requires: \usepackage{graphicx}

% ─── Phase 1 figures ─────────────────────────────────────────────────

\begin{figure}[!t]
  \centering
  \includegraphics[width=\columnwidth]{figures/F4_transmission_amplitude}
  \caption{T-RIS PIN-diode transmission amplitude $|T_n|$ vs frequency. On-state ($V_f > 0$) and off-state ($V_r < 0$) responses at 28 GHz carrier (vertical dashed line).}
  \label{fig:pin_transmission}
\end{figure}

\begin{figure}[!t]
  \centering
  \includegraphics[width=\columnwidth]{figures/F5_phase_response}
  \caption{T-RIS PIN-diode phase response $\angle T_n(f)$ in the on and off states. Phase difference $\Delta\phi = 64.5^\circ$ at 28 GHz.}
  \label{fig:pin_phase}
\end{figure}

\begin{figure}[!t]
  \centering
  \includegraphics[width=\columnwidth]{figures/F6_impedance_magnitude}
  \caption{PIN-diode impedance magnitude $|Z_n(f)|$ for forward and reverse bias.}
  \label{fig:pin_impedance}
\end{figure}

% ─── Phase 2 figures ─────────────────────────────────────────────────

\begin{figure}[!t]
  \centering
  \includegraphics[width=\columnwidth]{figures/F7_mfsk_spectrum}
  \caption{M-FSK tone spectrum at T-RIS output. Twelve sidebands at $f_c + k\Delta f$, $k=1,\ldots,12$, within the BPF passband.}
  \label{fig:mfsk_spectrum}
\end{figure}

\begin{figure*}[!t]
  \centering
  \includegraphics[width=\textwidth]{figures/F8_mfsk_spectrogram}
  \caption{Time--frequency spectrogram of the T-RIS M-FSK waveform. Each row shows one symbol period $T_{\rm sym} = 1.40\ \mu$s.}
  \label{fig:mfsk_spectrogram}
\end{figure*}

% ─── Phase 3 figures ─────────────────────────────────────────────────

\begin{figure}[!t]
  \centering
  \includegraphics[width=\columnwidth]{figures/F10_link_budget_waterfall}
  \caption{Bistatic link budget waterfall. Per-stage contributions (bars) and cumulative signal power (dashed) across the four-hop Tx $\to$ T-RIS $\to$ target $\to$ R-RIS $\to$ Rx path.}
  \label{fig:link_budget}
\end{figure}

\begin{figure*}[!t]
  \centering
  \includegraphics[width=\textwidth]{figures/F11_range_doppler}
  \caption{(a) IFFT range profile from $N_{\rm sw}=1024$ sweeps: target peak at 9.99 m (true 10 m, error $<1$ cm). (b) Range--Doppler map: target at $(10.0\ {\rm m},\ 2.5\ {\rm m/s})$.}
  \label{fig:range_doppler}
\end{figure*}

% ─── Phase 4 figures ─────────────────────────────────────────────────

\begin{figure*}[!t]
  \centering
  \includegraphics[width=\textwidth]{figures/F12_rmse_range_position}
  \caption{(a) Range RMSE and CRB vs single-symbol SNR. (b) 2-D position RMSE and CRB. Vertical dotted line: operating point SNR$_1 = -14.5$ dB.}
  \label{fig:rmse_range_pos}
\end{figure*}

\begin{figure}[!t]
  \centering
  \includegraphics[width=\columnwidth]{figures/F13_rmse_velocity}
  \caption{Velocity RMSE, CRB, and FFT estimation floor vs SNR. Three regimes: noise-dominated, CRB-tracking, and FFT-bin-limited.}
  \label{fig:rmse_velocity}
\end{figure}

% ─── Phase 5 figures ─────────────────────────────────────────────────

\begin{figure*}[!t]
  \centering
  \includegraphics[width=\textwidth]{figures/F14_tracking_trajectory}
  \caption{(a) Closed-loop beam-tracking scene over 50 CPIs ($T_{\rm tot}=860$ ms). Beam-steering arrows and colour-coded estimates. (b) Angle commands: true $\theta_{\rm true}$, noisy $\hat{\theta}$, and $\alpha$-smoothed $\tilde{\theta}$ ($\alpha=0.3$) vs CPI frame.}
  \label{fig:tracking_traj}
\end{figure*}

\begin{figure*}[!t]
  \centering
  \includegraphics[width=\textwidth]{figures/F15_tracking_errors}
  \caption{(a) Position-error scatter: 50th-percentile circle $= 5.3$ cm, 95th $= 9.2$ cm. (b) Empirical CDFs of position, range, and tangential errors; CRB reference $\sigma_R = 1.32$ cm.}
  \label{fig:tracking_err}
\end{figure*}

% ─── Phase 6 figures ─────────────────────────────────────────────────

\begin{figure*}[!t]
  \centering
  \includegraphics[width=\textwidth]{figures/F16_multitarget_vs_fmcw}
  \caption{Multi-target range profiles. (a) M-FSK IFFT: three clean peaks at true ranges 4, 7, 10 m regardless of target velocity. (b) FMCW FFT (same $B$, same CPI): Doppler-induced displacements of $+4.82$ m and $-3.21$ m create ghost targets.}
  \label{fig:multitarget}
\end{figure*}

\begin{figure}[!t]
  \centering
  \includegraphics[width=\columnwidth]{figures/F17_power_comparison}
  \caption{Stacked power breakdown for four sensing architectures. The two-stage T-RIS/R-RIS system achieves $2.33$ W, a $65\%$ reduction vs AESA ($6.56$ W).}
  \label{fig:power_bar}
\end{figure}

\begin{figure}[!t]
  \centering
  \includegraphics[width=\columnwidth]{figures/F18_spider_chart}
  \caption{Normalised performance spider chart (five axes). The proposed two-stage architecture (solid blue) leads in symbol rate and beam agility while matching the baseline in range resolution and hardware cost.}
  \label{fig:spider}
\end{figure}

% ─── Phase 7 figures ─────────────────────────────────────────────────

\begin{figure*}[!t]
  \centering
  \includegraphics[width=\textwidth]{figures/FS1_element_tolerance}
  \caption{Element-level manufacturing tolerance robustness. (a) Beam-gain change vs amplitude error $\sigma_A$: negligible impact ($<0.001$ dB at $\sigma_A=2\%$). (b) Beam-gain loss vs phase error $\sigma_\phi$: $0.012$ dB at nominal $3^\circ$ tolerance.}
  \label{fig:tolerance}
\end{figure*}

\begin{figure*}[!t]
  \centering
  \includegraphics[width=\textwidth]{figures/FS2_gap_and_quantisation}
  \caption{Hardware sensitivity. (a) SNR$_{\rm CPI}$ and CRB$_R$ vs inter-stage gap $d_g$: $5.10$ dB/mm sensitivity drives $\pm 0.196$ mm assembly tolerance. (b) Beam-gain loss vs phase bits $B$: $B=3$ bits incurs only $0.22$ dB loss.}
  \label{fig:gap_quant}
\end{figure*}
